\begin{cabstract}
	\addcontentsline{toc}{chapter}{摘要}
	3/2 随机波动率模型凭借其非线性扩散结构，在刻画资产价格极端波动、拟合波动率微笑形态上具备显著的理论优势，是高波动市场衍生品定价的核心工具之一。然而现有研究对该模型下各类定价方法的性能评价，多局限于平稳市场的理想化学术基准算例，缺乏真实极端市场环境下的实证检验，不同方法的工程适用性边界尚未得到系统性厘清。
	针对这一问题，本文在纯扩散框架下首次将Poisson‑Gamma 近似精确步进、拟指数（Quasi‑Exponential, QE）近似精确步进、逆高斯（Inverse Gaussian, IG）近似精确模拟、Gamma 近似精确模拟以及傅里叶余弦（COS）方法系统应用于 3/2 模型欧式期权定价，构建起包含时间步进类条件蒙特卡洛、终端精确与近似精确模拟、傅里叶频域定价三大范式、共九类方法的完整理论体系，并完整推导了各类方法的实现逻辑、误差特性与收敛性质。在此基础上，本文通过 7 组覆盖常规市场与极端参数环境的数值实验，量化对比了不同方法的定价精度、计算效率与全场景数值稳定性，形成了清晰的方法性能分层体系。
	进一步地，本文以 2025 年 10 月比特币市场闪崩极端事件为研究对象，构建 “长期固定参数时间序列估计 - 时变参数横截面校准 - 定价方法极端场景再检验” 的完整实证框架，验证了 3/2 模型对加密货币极端行情的市场适配性，识别了崩盘事件全周期模型核心参数的事件驱动型时变规律，并完成了定价方法在真实极端市场环境下的性能再检验。
	研究结果表明：傅里叶余弦（Fourier-Cos）方法是欧式期权定价的全局最优方案，在全场景兼具极致的定价精度、数值稳定性与计算效率；Poisson-Gamma 与 Quasi-Exponential（QE）近似精确步进方法，是蒙特卡洛框架下唯一全场景无失效的方案，为路径依赖型衍生品定价与动态对冲提供了可靠实现路径；标称 “精确” 的终端模拟方法仅能在理想化基准参数下完成理论验证，在极端市场环境中普遍出现数值溢出与系统性定价偏差，不具备工程实践层面的适配性。
	本文的研究弥补了现有 3/2 模型定价研究中理论算例与真实市场脱节的不足，形成了不同应用场景下的方法选择指引，为高波动市场环境下的衍生品定价、模型校准与风险管理提供了系统的理论支撑与实证依据。
\end{cabstract}

% vim:ts=4:sw=4